
\chapter{Introduction}

\section{Hydrogenic atoms}
\begin{figure}[htbp]
    \centering
    \includegraphics[width=1.0\linewidth]{Figures/Chapter1_Panel.png}
    \caption{Caption}
    \label{fig:placeholder}
\end{figure}
Hydrogenic atoms are atomic systems containing a single
electron bound to a nucleus of charge \(+Ze\), where
\(Z\) is the atomic number. The simplest example is
hydrogen (\(Z=1\)), although the same description applies
to one-electron ions such as \(\mathrm{He}^{+}\),
\(\mathrm{Li}^{2+}\), and \(\mathrm{Mn}^{24+}\).

Unlike multi-electron atoms, hydrogenic systems contain
no electron--electron interactions. Their nonrelativistic
Schrödinger equation can therefore be solved analytically,
providing exact wavefunctions and energy eigenvalues
against which numerical methods can be tested. \textit{(Insert Bransden and Joachain reference upon revision)} ...

In atomic units, and assuming an infinitely massive
point nucleus, the Hamiltonian is

\begin{equation}
    \hat{H}
    = -\frac{1}{2}\nabla^2-\frac{Z}{r}
    \label{eq:intro_hamiltonian}
\end{equation}

where \(r\) is the electron--nucleus separation.
The corresponding ground-state energy and normalised
\(1s\) wavefunction are

\begin{equation}
    E_{1s}=-\frac{Z^2}{2},
    \qquad
    \psi_{1s}(r)
    =\sqrt{\frac{Z^3}{\pi}}e^{-Zr}
    \label{eq:intro_hydrogenic}
\end{equation}

The exponential dependence describes the spatial decay
of the electron wavefunction away from the nucleus.
As \(Z\) increases, the electron becomes more tightly
bound and its characteristic spatial extent decreases.

Hydrogenic systems are consequently useful benchmarks
for investigating the accuracy, convergence and numerical
stability of approximate quantum-mechanical methods.


\section{Gaussian basis sets}

Although the hydrogenic Schrödinger equation has an
analytical solution, most atomic and molecular systems
cannot be solved exactly. Numerical electronic-structure
methods therefore commonly approximate wavefunctions
using finite expansions of simpler basis functions.

One widely used approach employs Gaussian-type orbitals
, which were introduced into electronic-structure
calculations by Boys \cite{Boys1950}.

A spherical Gaussian basis function has the form

\begin{equation}
    \phi_i(r)=A_i e^{-\lambda_i r^2}
    \label{eq:intro_gaussian}
\end{equation}

where \(\lambda_i>0\) determines the spatial width
and \(A_i\) is a normalisation constant.

Small exponents produce broad functions extending
farther from the nucleus, whereas large exponents
produce narrow functions concentrated near the nucleus.

The hydrogenic ground-state wavefunction can be
approximated by a linear combination of \(N\)
Gaussian functions,

\begin{equation}
    \psi_{\mathrm G}(r)
    =\sum_{i=1}^{N}c_i\phi_i(r)
    \label{eq:intro_expansion}
\end{equation}

where the coefficients \(c_i\) determine the
contribution of each basis function.

The coefficients are obtained by minimising the
energy expectation value, leading to a generalised
matrix eigenvalue problem. The resulting eigenvalues
provide approximate energies, while the eigenvectors
determine the corresponding wavefunctions.

The principal computational advantage of Gaussian
functions is that their products and many associated
integrals can be evaluated analytically. This makes
them particularly convenient for constructing the
Hamiltonian and overlap matrices required in
electronic-structure calculations \cite{Boys1950}.

Gaussian basis sets have consequently become a
standard component of quantum-chemical calculations,
including more elaborate contracted and correlation-
consistent basis sets \cite{Dunning1989,Raffenetti}.


\section{Accuracy and limitations of Gaussian expansions}

Despite their computational advantages, Gaussian
functions do not have the same radial dependence
as exact hydrogenic orbitals.

The exact \(1s\) wavefunction decays as \(e^{-Zr}\),
whereas an individual Gaussian decays as
\(e^{-\lambda r^2}\). A finite combination of
Gaussians can approximate the exponential function
closely over a range of distances, but cannot
reproduce its behaviour exactly everywhere.

An important difference occurs at the nucleus.
The exact hydrogenic wavefunction satisfies the
electron--nucleus cusp condition,

\begin{equation}
    \left.
    \frac{\mathrm d\psi_{1s}}{\mathrm dr}
    \right|_{r=0}
    =-Z\psi_{1s}(0)
    \label{eq:intro_cusp}
\end{equation}

whereas the radial derivative of any finite
combination of centred spherical Gaussians
vanishes at \(r=0\).

This limitation is connected to the Coulomb
singularity and the behaviour of electronic
wavefunctions near nuclei \cite{Kato}.
The theoretical accuracy of Gaussian expansions
has also been investigated by Kutzelnigg
\cite{Kutzelnigg1994}.

Increasing the number of basis functions generally
allows a more flexible approximation, although
the rate of improvement depends on the chosen
Gaussian exponents.

In particular, a basis containing functions
with very similar widths may become nearly
linearly dependent, leading to an ill-conditioned
overlap matrix and potential numerical instability
\cite{Lehtola2019}.

The selection of Gaussian exponents is therefore
as important as the number of basis functions.


\section{Even-tempered Gaussian basis sets}
Here the term \textit{even-tempered} refers to the constant
ratio between successive Gaussian exponents, i.e.
$\lambda_{i+1}/\lambda_i=\beta$, meaning that
the exponents are equally spaced on a logarithmic scale.
A systematic way of selecting Gaussian exponents
is to use a geometric progression,

\begin{equation}
    \lambda_i=\alpha\beta^{i-1},
    \qquad i=1,\ldots,N
    \label{eq:intro_even_tempered}
\end{equation}

where \(\alpha>0\) sets the smallest exponent,
\(\beta>1\) controls the spacing between successive
exponents, and \(N\) is the number of functions.

These are commonly referred to as even-tempered
Gaussian basis sets. Their systematic construction
makes it possible to investigate the influence of
basis size and exponent distribution independently.

Even-tempered bases and the optimisation of Gaussian
exponents have been studied extensively by Feller
and Ruedenberg \cite{FellerRuedenberg1979},
Bardo and Ruedenberg \cite{BaR},
and Wilson \cite{Wilson}.

For the present investigation, the parameters
\(N\), \(\alpha\), and \(\beta\) provide a
convenient means of controlling the flexibility
and numerical behaviour of the basis.


\section{Motivation and scope}

The hydrogenic ground state provides a useful
test case because both its energy and wavefunction
are known analytically. This allows numerical
accuracy to be assessed without uncertainty
arising from the reference solution.

The present investigation considers not only
the convergence of the ground-state energy,
but also the accuracy of the reconstructed
wavefunction and radial probability distribution.

This distinction is important because a small
energy error does not necessarily imply an
equally small error in every observable or
in the local behaviour of the wavefunction
\cite{Kutzelnigg1994}.

A further question concerns the transferability
of a Gaussian basis between hydrogenic systems.
Increasing nuclear charge contracts the exact
wavefunction, suggesting that a basis optimised
for hydrogen may be unsuitable for more
strongly bound one-electron ions.

The dependence of hydrogenic length scales on
nuclear charge provides a natural way to
investigate whether Gaussian exponents can
be rescaled to preserve accuracy.


\section{Aims and objectives}

The aim of this report is to investigate the
accuracy, convergence and transferability of
finite Gaussian basis approximations to
hydrogenic ground states.

The principal objectives are:

\begin{enumerate}
    \item Construct a Gaussian basis approximation
    to the hydrogen \(1s\) wavefunction and
    compare it with the exact solution.

    \item Investigate how the ground-state energy
    error depends on the number of Gaussian
    basis functions.

    \item Examine the influence of the exponent
    parameters \(\alpha\) and \(\beta\).

    \item Assess numerical stability through
    the condition number of the overlap matrix.

    \item Compare energy accuracy with
    wavefunction and radial-observable accuracy.

    \item Investigate the transferability of
    a hydrogen basis to higher Z hydrogenic ions,
    up to and including \(\mathrm{Mn}^{24+}\).

    \item Examine whether nuclear-charge
    scaling of the Gaussian exponents
    preserves the accuracy of the approximation.
\end{enumerate}


\section{Report structure}

\begin{description}
    \item[\textbf{Chapter 2:}] Introduces the theoretical framework,
    including Gaussian normalisation, the variational method,
    matrix elements and accuracy measures.

    \item[\textbf{Chapter 3:}] Describes the numerical implementation
    and the parameter investigations.

    \item[\textbf{Chapter 4:}] Presents the results for hydrogen,
    basis-set convergence, exponent selection,
    numerical conditioning and nuclear-charge scaling.

    \item[\textbf{Chapter 5:}] Summarises the main findings,
    limitations and possible extensions.
\end{description}


